% Options for packages loaded elsewhere
\PassOptionsToPackage{unicode}{hyperref}
\PassOptionsToPackage{hyphens}{url}
%
\documentclass[
  ignorenonframetext,
  aspectratio=169,
]{beamer}
\usepackage{pgfpages}
\setbeamertemplate{caption}[numbered]
\setbeamertemplate{caption label separator}{: }
\setbeamercolor{caption name}{fg=normal text.fg}
\beamertemplatenavigationsymbolshorizontal
% Prevent slide breaks in the middle of a paragraph
\widowpenalties 1 10000
\raggedbottom
\setbeamertemplate{part page}{
  \centering
  \begin{beamercolorbox}[sep=16pt,center]{part title}
    \usebeamerfont{part title}\insertpart\par
  \end{beamercolorbox}
}
\setbeamertemplate{section page}{
  \centering
  \begin{beamercolorbox}[sep=12pt,center]{part title}
    \usebeamerfont{section title}\insertsection\par
  \end{beamercolorbox}
}
\setbeamertemplate{subsection page}{
  \centering
  \begin{beamercolorbox}[sep=8pt,center]{part title}
    \usebeamerfont{subsection title}\insertsubsection\par
  \end{beamercolorbox}
}
\AtBeginPart{
  \frame{\partpage}
}
\AtBeginSection{
  \ifbibliography
  \else
    \frame{\sectionpage}
  \fi
}
\AtBeginSubsection{
  \frame{\subsectionpage}
}

\usepackage{amsmath,amssymb}
\usepackage{iftex}
\ifPDFTeX
  \usepackage[T1]{fontenc}
  \usepackage[utf8]{inputenc}
  \usepackage{textcomp} % provide euro and other symbols
\else % if luatex or xetex
  \usepackage{unicode-math}
  \defaultfontfeatures{Scale=MatchLowercase}
  \defaultfontfeatures[\rmfamily]{Ligatures=TeX,Scale=1}
\fi
\usepackage{lmodern}
\usetheme[]{Montpellier}
\usecolortheme{seagull}
\ifPDFTeX\else  
    % xetex/luatex font selection
\fi
% Use upquote if available, for straight quotes in verbatim environments
\IfFileExists{upquote.sty}{\usepackage{upquote}}{}
\IfFileExists{microtype.sty}{% use microtype if available
  \usepackage[]{microtype}
  \UseMicrotypeSet[protrusion]{basicmath} % disable protrusion for tt fonts
}{}
\usepackage{xcolor}
\newif\ifbibliography
\setlength{\emergencystretch}{3em} % prevent overfull lines
\setcounter{secnumdepth}{5}


\providecommand{\tightlist}{%
  \setlength{\itemsep}{0pt}\setlength{\parskip}{0pt}}\usepackage{longtable,booktabs,array}
\usepackage{calc} % for calculating minipage widths
\usepackage{caption}
% Make caption package work with longtable
\makeatletter
\def\fnum@table{\tablename~\thetable}
\makeatother
\usepackage{graphicx}
\makeatletter
\def\maxwidth{\ifdim\Gin@nat@width>\linewidth\linewidth\else\Gin@nat@width\fi}
\def\maxheight{\ifdim\Gin@nat@height>\textheight\textheight\else\Gin@nat@height\fi}
\makeatother
% Scale images if necessary, so that they will not overflow the page
% margins by default, and it is still possible to overwrite the defaults
% using explicit options in \includegraphics[width, height, ...]{}
\setkeys{Gin}{width=\maxwidth,height=\maxheight,keepaspectratio}
% Set default figure placement to htbp
\makeatletter
\def\fps@figure{htbp}
\makeatother

\usepackage{libertine}
\makeatletter
\makeatother
\makeatletter
\makeatother
\makeatletter
\@ifpackageloaded{caption}{}{\usepackage{caption}}
\AtBeginDocument{%
\ifdefined\contentsname
  \renewcommand*\contentsname{Содержание}
\else
  \newcommand\contentsname{Содержание}
\fi
\ifdefined\listfigurename
  \renewcommand*\listfigurename{Список иллюстраций}
\else
  \newcommand\listfigurename{Список иллюстраций}
\fi
\ifdefined\listtablename
  \renewcommand*\listtablename{Список таблиц}
\else
  \newcommand\listtablename{Список таблиц}
\fi
\ifdefined\figurename
  \renewcommand*\figurename{Рисунок}
\else
  \newcommand\figurename{Рисунок}
\fi
\ifdefined\tablename
  \renewcommand*\tablename{Таблица}
\else
  \newcommand\tablename{Таблица}
\fi
}
\@ifpackageloaded{float}{}{\usepackage{float}}
\floatstyle{ruled}
\@ifundefined{c@chapter}{\newfloat{codelisting}{h}{lop}}{\newfloat{codelisting}{h}{lop}[chapter]}
\floatname{codelisting}{Список}
\newcommand*\listoflistings{\listof{codelisting}{Листинги}}
\makeatother
\makeatletter
\@ifpackageloaded{caption}{}{\usepackage{caption}}
\@ifpackageloaded{subcaption}{}{\usepackage{subcaption}}
\makeatother
\makeatletter
\@ifpackageloaded{tcolorbox}{}{\usepackage[skins,breakable]{tcolorbox}}
\makeatother
\makeatletter
\@ifundefined{shadecolor}{\definecolor{shadecolor}{rgb}{.97, .97, .97}}
\makeatother
\makeatletter
\makeatother
\makeatletter
\makeatother
\ifLuaTeX
\usepackage[bidi=basic]{babel}
\else
\usepackage[bidi=default]{babel}
\fi
\babelprovide[main,import]{russian}
\babelprovide[import]{english}
% get rid of language-specific shorthands (see #6817):
\let\LanguageShortHands\languageshorthands
\def\languageshorthands#1{}
\ifLuaTeX
  \usepackage{selnolig}  % disable illegal ligatures
\fi
\usepackage{csquotes}
\IfFileExists{bookmark.sty}{\usepackage{bookmark}}{\usepackage{hyperref}}
\IfFileExists{xurl.sty}{\usepackage{xurl}}{} % add URL line breaks if available
\urlstyle{same} % disable monospaced font for URLs
\hypersetup{
  pdftitle={Лабораторная работа № 5},
  pdfauthor={Оразгелдиев Язгелди},
  pdflang={ru-RU},
  hidelinks,
  pdfcreator={LaTeX via pandoc}}

\title{Лабораторная работа № 5}
\subtitle{Построение графиков}
\author{Оразгелдиев Язгелди}
\date{2025-11-22}

\begin{document}
\frame{\titlepage}
\ifdefined\Shaded\renewenvironment{Shaded}{\begin{tcolorbox}[interior hidden, frame hidden, boxrule=0pt, sharp corners, borderline west={3pt}{0pt}{shadecolor}, enhanced, breakable]}{\end{tcolorbox}}\fi

\renewcommand*\contentsname{Содержание}
\begin{frame}[allowframebreaks]
  \frametitle{Содержание}
  \tableofcontents[hideallsubsections]
\end{frame}
\hypertarget{ux438ux43dux444ux43eux440ux43cux430ux446ux438ux44f}{%
\section{Информация}\label{ux438ux43dux444ux43eux440ux43cux430ux446ux438ux44f}}

\begin{frame}{Докладчик}
\protect\hypertarget{ux434ux43eux43aux43bux430ux434ux447ux438ux43a}{}
\begin{itemize}[<+->]
\tightlist
\item
  Оразгелдиев Язгелди
\item
  \href{mailto:1032225075@pfur.ru}{\nolinkurl{1032225075@pfur.ru}}
\item
  \url{https://github.com/YazgeldiOrazgeldiyev}
\end{itemize}
\end{frame}

\begin{frame}{Цель работы}
\protect\hypertarget{ux446ux435ux43bux44c-ux440ux430ux431ux43eux442ux44b}{}
Основная цель работы -- освоить синтаксис языка Julia для построения
графиков.
\end{frame}

\begin{frame}{Задание}
\protect\hypertarget{ux437ux430ux434ux430ux43dux438ux435}{}
\begin{enumerate}[<+->]
\item
  Используя JupyterLab, повторите примерыи. При этом дополните графики
  обозначениями осей координат, легендой с названиями траекторий,
  названиями графиков и т.п.
\item
  Выполните задания для самостоятельной работы.
\end{enumerate}
\end{frame}

\begin{frame}{Теоретическое введение}
\protect\hypertarget{ux442ux435ux43eux440ux435ux442ux438ux447ux435ux441ux43aux43eux435-ux432ux432ux435ux434ux435ux43dux438ux435}{}
Julia -- высокоуровневый свободный язык программирования с динамической
типизацией, созданный для математических вычислений {[}@julialang{]}.
Эффективен также и для написания программ общего назначения. Синтаксис
языка схож с синтаксисом других математических языков, однако имеет
некоторые существенные отличия.

Для выполнения заданий была использована официальная документация Julia
{[}@juliadoc{]}.
\end{frame}

\begin{frame}{Выполнение лабораторной работы}
\protect\hypertarget{ux432ux44bux43fux43eux43bux43dux435ux43dux438ux435-ux43bux430ux431ux43eux440ux430ux442ux43eux440ux43dux43eux439-ux440ux430ux431ux43eux442ux44b}{}
Выполним примеры из лабораторной работы для знакомства с пакетами по
отрисовки графиков и их функциями (рис.
{[}-@fig:001{]}-{[}-@fig:019{]}).

\begin{figure}

{\centering \includegraphics[width=0.25\textwidth,height=\textheight]{image/1.png}

}

\caption{Основные пакеты для работы с графиками в Julia}

\end{figure}
\end{frame}

\begin{frame}{Выполнение лабораторной работы}
\protect\hypertarget{ux432ux44bux43fux43eux43bux43dux435ux43dux438ux435-ux43bux430ux431ux43eux440ux430ux442ux43eux440ux43dux43eux439-ux440ux430ux431ux43eux442ux44b-1}{}
\begin{figure}

{\centering \includegraphics[width=0.25\textwidth,height=\textheight]{image/2.png}

}

\caption{Основные пакеты для работы с графиками в Julia}

\end{figure}
\end{frame}

\begin{frame}{Выполнение лабораторной работы}
\protect\hypertarget{ux432ux44bux43fux43eux43bux43dux435ux43dux438ux435-ux43bux430ux431ux43eux440ux430ux442ux43eux440ux43dux43eux439-ux440ux430ux431ux43eux442ux44b-2}{}
\begin{figure}

{\centering \includegraphics[width=0.25\textwidth,height=\textheight]{image/3.png}

}

\caption{Опции при построении графика}

\end{figure}
\end{frame}

\begin{frame}{Выполнение лабораторной работы}
\protect\hypertarget{ux432ux44bux43fux43eux43bux43dux435ux43dux438ux435-ux43bux430ux431ux43eux440ux430ux442ux43eux440ux43dux43eux439-ux440ux430ux431ux43eux442ux44b-3}{}
\begin{figure}

{\centering \includegraphics[width=0.25\textwidth,height=\textheight]{image/4.png}

}

\caption{Опции при построении графика}

\end{figure}
\end{frame}

\begin{frame}{Выполнение лабораторной работы}
\protect\hypertarget{ux432ux44bux43fux43eux43bux43dux435ux43dux438ux435-ux43bux430ux431ux43eux440ux430ux442ux43eux440ux43dux43eux439-ux440ux430ux431ux43eux442ux44b-4}{}
\begin{figure}

{\centering \includegraphics[width=0.25\textwidth,height=\textheight]{image/5.png}

}

\caption{Точечный график}

\end{figure}
\end{frame}

\begin{frame}{Выполнение лабораторной работы}
\protect\hypertarget{ux432ux44bux43fux43eux43bux43dux435ux43dux438ux435-ux43bux430ux431ux43eux440ux430ux442ux43eux440ux43dux43eux439-ux440ux430ux431ux43eux442ux44b-5}{}
\begin{figure}

{\centering \includegraphics[width=0.25\textwidth,height=\textheight]{image/6.png}

}

\caption{Точечный график}

\end{figure}
\end{frame}

\begin{frame}{Выполнение лабораторной работы}
\protect\hypertarget{ux432ux44bux43fux43eux43bux43dux435ux43dux438ux435-ux43bux430ux431ux43eux440ux430ux442ux43eux440ux43dux43eux439-ux440ux430ux431ux43eux442ux44b-6}{}
\begin{figure}

{\centering \includegraphics[width=0.25\textwidth,height=\textheight]{image/7.png}

}

\caption{Аппроксимация данных}

\end{figure}
\end{frame}

\begin{frame}{Выполнение лабораторной работы}
\protect\hypertarget{ux432ux44bux43fux43eux43bux43dux435ux43dux438ux435-ux43bux430ux431ux43eux440ux430ux442ux43eux440ux43dux43eux439-ux440ux430ux431ux43eux442ux44b-7}{}
\begin{figure}

{\centering \includegraphics[width=0.25\textwidth,height=\textheight]{image/8.png}

}

\caption{Полярные координаты}

\end{figure}
\end{frame}

\begin{frame}{Выполнение лабораторной работы}
\protect\hypertarget{ux432ux44bux43fux43eux43bux43dux435ux43dux438ux435-ux43bux430ux431ux43eux440ux430ux442ux43eux440ux43dux43eux439-ux440ux430ux431ux43eux442ux44b-8}{}
\begin{figure}

{\centering \includegraphics[width=0.25\textwidth,height=\textheight]{image/9.png}

}

\caption{График поверхности}

\end{figure}
\end{frame}

\begin{frame}{Выполнение лабораторной работы}
\protect\hypertarget{ux432ux44bux43fux43eux43bux43dux435ux43dux438ux435-ux43bux430ux431ux43eux440ux430ux442ux43eux440ux43dux43eux439-ux440ux430ux431ux43eux442ux44b-9}{}
\begin{figure}

{\centering \includegraphics[width=0.25\textwidth,height=\textheight]{image/10.png}

}

\caption{Линии уровня}

\end{figure}
\end{frame}

\begin{frame}{Выполнение лабораторной работы}
\protect\hypertarget{ux432ux44bux43fux43eux43bux43dux435ux43dux438ux435-ux43bux430ux431ux43eux440ux430ux442ux43eux440ux43dux43eux439-ux440ux430ux431ux43eux442ux44b-10}{}
\begin{figure}

{\centering \includegraphics[width=0.25\textwidth,height=\textheight]{image/11.png}

}

\caption{Векторные поля}

\end{figure}
\end{frame}

\begin{frame}{Выполнение лабораторной работы}
\protect\hypertarget{ux432ux44bux43fux43eux43bux43dux435ux43dux438ux435-ux43bux430ux431ux43eux440ux430ux442ux43eux440ux43dux43eux439-ux440ux430ux431ux43eux442ux44b-11}{}
\begin{figure}

{\centering \includegraphics[width=0.25\textwidth,height=\textheight]{image/12.png}

}

\caption{Анимация}

\end{figure}
\end{frame}

\begin{frame}{Выполнение лабораторной работы}
\protect\hypertarget{ux432ux44bux43fux43eux43bux43dux435ux43dux438ux435-ux43bux430ux431ux43eux440ux430ux442ux43eux440ux43dux43eux439-ux440ux430ux431ux43eux442ux44b-12}{}
\begin{figure}

{\centering \includegraphics[width=0.25\textwidth,height=\textheight]{image/13.png}

}

\caption{Гипоциклоида}

\end{figure}
\end{frame}

\begin{frame}{Выполнение лабораторной работы}
\protect\hypertarget{ux432ux44bux43fux43eux43bux43dux435ux43dux438ux435-ux43bux430ux431ux43eux440ux430ux442ux43eux440ux43dux43eux439-ux440ux430ux431ux43eux442ux44b-13}{}
\begin{figure}

{\centering \includegraphics[width=0.25\textwidth,height=\textheight]{image/14.png}

}

\caption{Errorbars}

\end{figure}
\end{frame}

\begin{frame}{Выполнение лабораторной работы}
\protect\hypertarget{ux432ux44bux43fux43eux43bux43dux435ux43dux438ux435-ux43bux430ux431ux43eux440ux430ux442ux43eux440ux43dux43eux439-ux440ux430ux431ux43eux442ux44b-14}{}
\begin{figure}

{\centering \includegraphics[width=0.25\textwidth,height=\textheight]{image/15.png}

}

\caption{Errorbars}

\end{figure}
\end{frame}

\begin{frame}{Выполнение лабораторной работы}
\protect\hypertarget{ux432ux44bux43fux43eux43bux43dux435ux43dux438ux435-ux43bux430ux431ux43eux440ux430ux442ux43eux440ux43dux43eux439-ux440ux430ux431ux43eux442ux44b-15}{}
\begin{figure}

{\centering \includegraphics[width=0.25\textwidth,height=\textheight]{image/16.png}

}

\caption{Использование пакета Distributions}

\end{figure}
\end{frame}

\begin{frame}{Выполнение лабораторной работы}
\protect\hypertarget{ux432ux44bux43fux43eux43bux43dux435ux43dux438ux435-ux43bux430ux431ux43eux440ux430ux442ux43eux440ux43dux43eux439-ux440ux430ux431ux43eux442ux44b-16}{}
\begin{figure}

{\centering \includegraphics[width=0.25\textwidth,height=\textheight]{image/17.png}

}

\caption{Подграфики}

\end{figure}
\end{frame}

\begin{frame}{Выполнение лабораторной работы}
\protect\hypertarget{ux432ux44bux43fux43eux43bux43dux435ux43dux438ux435-ux43bux430ux431ux43eux440ux430ux442ux43eux440ux43dux43eux439-ux440ux430ux431ux43eux442ux44b-17}{}
\begin{figure}

{\centering \includegraphics[width=0.25\textwidth,height=\textheight]{image/18.png}

}

\caption{Подграфики}

\end{figure}
\end{frame}

\begin{frame}{Выполнение лабораторной работы}
\protect\hypertarget{ux432ux44bux43fux43eux43bux43dux435ux43dux438ux435-ux43bux430ux431ux43eux440ux430ux442ux43eux440ux43dux43eux439-ux440ux430ux431ux43eux442ux44b-18}{}
\begin{figure}

{\centering \includegraphics[width=0.25\textwidth,height=\textheight]{image/19.png}

}

\caption{Подграфики}

\end{figure}
\end{frame}

\begin{frame}{Задания для самостоятельного выполнения}
\protect\hypertarget{ux437ux430ux434ux430ux43dux438ux44f-ux434ux43bux44f-ux441ux430ux43cux43eux441ux442ux43eux44fux442ux435ux43bux44cux43dux43eux433ux43e-ux432ux44bux43fux43eux43bux43dux435ux43dux438ux44f}{}
\begin{figure}

{\centering \includegraphics[width=0.25\textwidth,height=\textheight]{image/20.png}

}

\caption{Задание №1 и №2}

\end{figure}
\end{frame}

\begin{frame}{Выполнение лабораторной работы}
\protect\hypertarget{ux432ux44bux43fux43eux43bux43dux435ux43dux438ux435-ux43bux430ux431ux43eux440ux430ux442ux43eux440ux43dux43eux439-ux440ux430ux431ux43eux442ux44b-19}{}
\begin{figure}

{\centering \includegraphics[width=0.25\textwidth,height=\textheight]{image/21.png}

}

\caption{Задание №2}

\end{figure}
\end{frame}

\begin{frame}{Выполнение лабораторной работы}
\protect\hypertarget{ux432ux44bux43fux43eux43bux43dux435ux43dux438ux435-ux43bux430ux431ux43eux440ux430ux442ux43eux440ux43dux43eux439-ux440ux430ux431ux43eux442ux44b-20}{}
\begin{figure}

{\centering \includegraphics[width=0.25\textwidth,height=\textheight]{image/22.png}

}

\caption{Задание №3}

\end{figure}
\end{frame}

\begin{frame}{Выполнение лабораторной работы}
\protect\hypertarget{ux432ux44bux43fux43eux43bux43dux435ux43dux438ux435-ux43bux430ux431ux43eux440ux430ux442ux43eux440ux43dux43eux439-ux440ux430ux431ux43eux442ux44b-21}{}
\begin{figure}

{\centering \includegraphics[width=0.25\textwidth,height=\textheight]{image/23.png}

}

\caption{Задание №4}

\end{figure}
\end{frame}

\begin{frame}{Выполнение лабораторной работы}
\protect\hypertarget{ux432ux44bux43fux43eux43bux43dux435ux43dux438ux435-ux43bux430ux431ux43eux440ux430ux442ux43eux440ux43dux43eux439-ux440ux430ux431ux43eux442ux44b-22}{}
\begin{figure}

{\centering \includegraphics[width=0.25\textwidth,height=\textheight]{image/24.png}

}

\caption{Задание №5}

\end{figure}
\end{frame}

\begin{frame}{Выполнение лабораторной работы}
\protect\hypertarget{ux432ux44bux43fux43eux43bux43dux435ux43dux438ux435-ux43bux430ux431ux43eux440ux430ux442ux43eux440ux43dux43eux439-ux440ux430ux431ux43eux442ux44b-23}{}
\begin{figure}

{\centering \includegraphics[width=0.25\textwidth,height=\textheight]{image/25.png}

}

\caption{Задание №6}

\end{figure}
\end{frame}

\begin{frame}{Выполнение лабораторной работы}
\protect\hypertarget{ux432ux44bux43fux43eux43bux43dux435ux43dux438ux435-ux43bux430ux431ux43eux440ux430ux442ux43eux440ux43dux43eux439-ux440ux430ux431ux43eux442ux44b-24}{}
\begin{figure}

{\centering \includegraphics[width=0.25\textwidth,height=\textheight]{image/26.png}

}

\caption{Задание №7}

\end{figure}
\end{frame}

\begin{frame}{Выполнение лабораторной работы}
\protect\hypertarget{ux432ux44bux43fux43eux43bux43dux435ux43dux438ux435-ux43bux430ux431ux43eux440ux430ux442ux43eux440ux43dux43eux439-ux440ux430ux431ux43eux442ux44b-25}{}
\begin{figure}

{\centering \includegraphics[width=0.25\textwidth,height=\textheight]{image/27.png}

}

\caption{Задание №8}

\end{figure}
\end{frame}

\begin{frame}{Выполнение лабораторной работы}
\protect\hypertarget{ux432ux44bux43fux43eux43bux43dux435ux43dux438ux435-ux43bux430ux431ux43eux440ux430ux442ux43eux440ux43dux43eux439-ux440ux430ux431ux43eux442ux44b-26}{}
\begin{figure}

{\centering \includegraphics[width=0.25\textwidth,height=\textheight]{image/28.png}

}

\caption{Задание №9}

\end{figure}
\end{frame}

\begin{frame}{Выполнение лабораторной работы}
\protect\hypertarget{ux432ux44bux43fux43eux43bux43dux435ux43dux438ux435-ux43bux430ux431ux43eux440ux430ux442ux43eux440ux43dux43eux439-ux440ux430ux431ux43eux442ux44b-27}{}
\begin{figure}

{\centering \includegraphics[width=0.25\textwidth,height=\textheight]{image/29.png}

}

\caption{Задание №10}

\end{figure}
\end{frame}

\begin{frame}{Выполнение лабораторной работы}
\protect\hypertarget{ux432ux44bux43fux43eux43bux43dux435ux43dux438ux435-ux43bux430ux431ux43eux440ux430ux442ux43eux440ux43dux43eux439-ux440ux430ux431ux43eux442ux44b-28}{}
\begin{figure}

{\centering \includegraphics[width=0.25\textwidth,height=\textheight]{image/30.png}

}

\caption{Задание №11}

\end{figure}
\end{frame}

\begin{frame}{Выводы}
\protect\hypertarget{ux432ux44bux432ux43eux434ux44b}{}
В результате выполнения данной лабораторной работы я освоил синтаксис
языка Julia для построения графиков.
\end{frame}



\end{document}
